\section{Experiments}
\label{sec:experiments}

This is a self-contained section that explains all experimental protocol.

\subsection{Paper-facing family map}

I think we can write the results in three main scientific families and have other details e.g. hyperparameter selection for each setting in the appendix.

\begin{table*}[ht!]
\centering
\footnotesize
\setlength{\tabcolsep}{4pt}
\begin{tabular}{p{0.18\textwidth}p{0.37\textwidth}p{0.35\textwidth}}
\toprule
Family & What belongs here & Suggestions \\
\midrule
\textbf{Cross-system forecasting} &
The full 29-system benchmark at the final 200k training budget, using the fixed sparse MLP anchor and the fixed LISTA comparator &
This is the benchmark headline. The message should be that the fair 29-system benchmark separates typical-system quality from breadth: the sparse MLP is the safer anchor through $H1000$, LISTA wins more systems, and LISTA becomes median-best only at the longest horizons. \\
\textbf{Hard-system forecasting} &
Smaller-step-size follow-up studies, long-horizon reevaluations of those same checkpoints, Kuramoto robustness and dimension checks, and matched parity/fairness controls &
Write this as one family about when smaller step size and latent structure help, and when they still do not. \\
\textbf{Basin-support and mechanism} &
Basin-support alignment, label-free clustering, the direct Kuramoto support audit, the corrected competitive Lotka-Volterra follow-up, and recurring-support local-linearity &
Write this as one family about when latent supports do and do not align with basins. \\
\textbf{Appendix-only provenance} & LISTA hyperparameter selection &
Mention only briefly in the main paper.\\
\bottomrule
\end{tabular}
\caption{\textbf{Paper-facing family map.}}
\label{tab:review-family-map-20260314}
\end{table*}

\subsection{Global reproduction protocol}

\paragraph{Software and deterministic conventions.}
Built-in systems and Dysts systems are both integrated with fourth-order
Runge-Kutta. Do not assume one universal seed set across the paper; the exact
seed coverage is part of each family definition below, and some benchmark
families intentionally use mixed seed counts across systems. All reported
tables use medians across seeds rather than best-seed results. The selected
checkpoint is always the checkpoint with the lowest \emph{final-step} validation
error on a fixed 200-step validation rollout from 16 held-out initial states. The rationale for choosing the final step instead of a horizon average is that errors tend to compound as a rollout goes longer.

Exact seed conventions are also fixed: a run with training seed $s$ uses $s$
as the master model seed, the fixed validation corpus is generated with seed
$s + 999999$, forecasting evaluation initial conditions are generated with seed
$s + 12345$, and support-structure analyses use evaluation seed 42 unless
explicitly stated otherwise.

\subsubsection{Model families.}

{\color{red}[UDAY: What is sparse about sparse MLP? From the description it would appear it is just multi-layer ReLU encoder with linear decoder.]}

\begin{table*}[ht]
\centering
\scriptsize
\setlength{\tabcolsep}{4pt}
\begin{tabular}{p{0.18\textwidth}p{0.23\textwidth}p{0.16\textwidth}p{0.33\textwidth}}
\toprule
Model family & Encoder and decoder & Latent dynamics & Main benchmark recipe \\
\midrule
\textbf{Sparse MLP anchor} &
Encoder with two hidden layers of width 64, bias terms enabled, and ReLU at the encoder output; one learned linear decoder with no hidden layer &
Dense latent transition &
200k steps, latent dimension 256, main learning rate $10^{-4}$, latent-transition learning rate $10^{-5}$, weight decay $10^{-4}$, alignment weight 1.0, reconstruction weight 0.03, prediction weight 1.0, sparsity weight 0.0025 \\
\textbf{Zero-sparsity MLP baseline} &
Same encoder and decoder as the sparse MLP anchor &
Dense latent transition &
Same as the sparse MLP anchor except the \textbf{sparsity weight is 0.0} \\
\textbf{LISTA comparator} &
Same 64-64 front end as the sparse MLP, followed by one LISTA refinement step after the initial shrinkage map; LISTA uses $\alpha = 0.15$ and effective threshold $\alpha/L$, with $L$ computed once from the initial decoder dictionary; normalized linear dictionary decoder &
Dense latent transition &
200k steps, latent dimension 256, main learning rate $5 \times 10^{-5}$, latent-transition learning rate $5 \times 10^{-6}$, weight decay $10^{-4}$, alignment weight 1.0, reconstruction weight 0.03, prediction weight 1.0, sparsity weight 0.003 \\
\textbf{Block-diagonal LISTA} &
Same encoder and decoder as standard LISTA &
Block-diagonal latent transition with block size 16 &
Same as standard LISTA except the sparsity weight is either 0.003 or 0.006, depending on the comparison \\
\textbf{Block-diagonal MLP} &
Same encoder and decoder as the sparse MLP anchor &
Block-diagonal latent transition with block size 16 &
Used only in fairness controls to isolate whether gains come from encoder structure or from block structure in the latent transition \\
\bottomrule
\end{tabular}
\caption{\textbf{Human-readable model families.} The sparse-MLP versus standard-LISTA comparison is budget-matched and front-end-matched, but not exact parameter-count matching because LISTA adds a recurrent encoder matrix.}
\label{tab:review-model-families-20260314}
\end{table*}

\paragraph{Exact encoder and decoder maps.}
The sparse MLP anchor uses
$h_1 = \mathrm{ReLU}(W_1 x + b_1)$,
$h_2 = \mathrm{ReLU}(W_2 h_1 + b_2)$,
$z = \mathrm{ReLU}(W_3 h_2 + b_3)$, so two hidden layers and one output layer for a total of three layers of weights. Then the prediction
$\hat{x} = W_{\mathrm{dec}} z$.
The dense and block-diagonal LISTA families use the same three-layer MLP to produce
$c$ instead of $z$, then apply
$u^{(0)} = \mathrm{shrink}(c, \alpha / L)$ and
$z = \mathrm{ReLU}(\mathrm{shrink}(u^{(0)} S + c, \alpha / L))$,
where $S$ is learned and
$L$ is 1.05 times a 10-step power-iteration estimate of
$\|W_d^\top W_d\|_2$.
The LISTA-family decoder learns a latent-to-state dictionary and renormalizes
each latent atom to unit norm at decode time before applying the linear decode.

\paragraph{Shared training protocol.}
Every paper-facing 200k training run uses 200000 optimization steps, mini-batch
size 256, latent dimension 256, and training windows of 9 consecutive states,
which provide 8 one-step transitions per example. Each step optimizes a weighted
sum of latent alignment loss, reconstruction loss, multi-step prediction loss,
and an L1 sparsity penalty on the predicted latent states. AdamW is used with
one learning rate for all non-transition parameters, a 10x smaller learning rate for
the latent Koopman linear transition, and zero weight decay on the latent transition.
% To reproduce the exact data stream rather than only the high-level recipe, use
% 8 persistent mini-batch random-number streams with seeds
% $s + 256 i$ for $i \in \{0,\dots,7\}$, where $s$ is the training seed, and use
% stream $t \bmod 8$ at optimization step $t$. 
Validation-based checkpoint selection is run every 500 training steps plus the final step. With horizon
$H = 8$ and observation dimension $d_x$, the paper-facing loss is
\[
\frac{1}{H}\left[
\lambda_{\mathrm{align}} \mathbb{E}\|z_{\mathrm{pred}} - z_{\mathrm{gt}}\|_2
+
 \frac{1}{\sqrt{d_x}}(\lambda_{\mathrm{rec}} \mathbb{E}\|\hat{x}_{\mathrm{rec}} - x_{\mathrm{gt}}\|_2
+
\lambda_{\mathrm{pred}} \mathbb{E}\|\hat{x}_{\mathrm{pred}} - x_{\mathrm{gt}}\|_2)
+
\lambda_{\mathrm{sparse}}  \mathbb{E}\|z_{\mathrm{pred}}\|_1\right].
\]
The reconstruction term is applied to decoded \emph{true} future latents, not
to the initial state, and the sparsity penalty is applied to the predicted
latent rollout. Checkpoint selection uses a 200-step every-step-reencoding validation rollout and chooses the checkpoint with the smallest final-step Euclidean validation error. Every-step-reencoding is a simpler proxy for validation-based checkpoint selection because we select the best reencoding period at inference time and do not have access to what the best reencoding period is at training. The ``best'' reencoding period also varies significantly between different systems, so keeping this proxy consistent is fair to all systems.

\paragraph{System-generation details for reproducibility.}
For Hopfield, the stored patterns are random $\pm 1$ patterns, and unless a
system seed is explicitly fixed in a targeted study, the stored patterns change
with the training seed, using generator seed $s + 29$. For competitive
Lotka-Volterra, unless a fixed interaction seed is explicitly stated, the
interaction matrix changes with the training seed; in the benchmark setting,
off-diagonal interaction coefficients are sampled uniformly on $[0,0.70]$,
symmetrized, and then the diagonal is set to 1.0. Without a fixed system seed,
the growth-rate generator uses seed $s + 44$ and the interaction-matrix
generator uses seed $s + 45$. The benchmark Kuramoto system uses identical
natural frequencies and therefore does not vary across seeds; random-frequency
variants use generator seed $s + 13$. The five canonical multiwell attractor
centers are fixed rather than resampled.

\paragraph{Dysts-specific rules.}
For Dysts systems \citep{gilpin_chaos_2021, gilpin_model_2023}, use Dysts 0.96 with the default parameterization and default initial condition of the named system. 
Standardize each state coordinate using the Dysts metadata mean and standard deviation. Reset states are sampled as the default initial condition plus Gaussian noise with standard deviation $0.2 \times$ the coordinate-wise Dysts standard deviation. 
Use a deterministic trajectory cache with 200 trajectories per split, 30000 steps per trajectory, 2000 warmup steps, and cache-build seeds 0, 1, and 2 for train, validation, and test respectively. 
The Dysts cache is independent of the model seed: changing the model seed changes model initialization and batch selection inside the cache but does not rebuild the train/validation/test cache contents.

\paragraph{Shared forecasting evaluation.}
Each selected checkpoint is evaluated on 100 initial conditions. Rollout modes
are no re-encoding, every-step re-encoding, and periodic re-encoding. Periodic
cadences are \{10, 25, 50, 100\} for built-in systems and
\{1, 5, 10, 20, 40, 60, 80, 100, 200, 300, 400, 500, 1000\} for Dysts systems.
For a horizon $H$, compute squared Euclidean prediction error at each forecast
step from 1 to $H$, average over those $H$ steps for each initial condition, and
then average over initial conditions. ``Best periodic re-encoding'' means the
lowest mean horizon error among the tested periodic cadences. Good systems satisfy
H1000 best-periodic error $< 10$, and catastrophic systems satisfy H1000
best-periodic error $\ge 1000$. This is an evaluation-time oracle over a fixed,
predeclared set of re-encoding cadences rather than a learned policy. Every
paper-facing forecasting table first takes the system-level median across the
available seeds for that model and system, then takes the cross-system median
across the 29 system-level medians; if a system has only 3 collected seeds in
the final benchmark family, those 3 seeds are used directly rather than
imputing missing runs. For corrected 4-basin competitive Lotka--Volterra, the
first 15-seed collector \texttt{collect\_200k\_15seed} was only an intermediate
mixed-coverage summary: \texttt{generic\_sparse\_blockdiag\_200k},
\texttt{lista\_blockdiag\_ns200k\_denseopt\_sc3em3}, and
\texttt{lista\_blockdiag\_ns200k\_denseopt\_sc6em3} already had 15 seeds, but
\texttt{generic\_sparse\_ns200k\_best} and
\texttt{lista\_dense\_promoted\_stage4} still had only 10. Any corrected-CLV
15-seed table in this document therefore uses the rebuilt
\texttt{collect\_200k\_15seed\_full} summary, which extends all five roots to
seeds 0--14.

\subsection{Benchmark system catalog}

\subsubsection{Built-in systems.}

{\color{red}{[UDAY: Since we have eight 2 dimensional dynamical system, it would be good to add the phase plots.]}}

{\color{red}{[UDAY: Being a devil's advocate:]
\begin{enumerate}
    \item Why is sampling rate $dt$ is different for the same dimensional systems ?
    \item Hopfield system seems to be sampled from Gaussian whereas the rest using uniform distribution.
    Suppose that Gaussian is $\mathcal{N}(0, \Sigma)$ what does standard deviation here means $\sqrt{\Sigma}$ or $\sqrt{\textrm{tr}(\Sigma)}$ ? Note that either term scales with square root of the dimension.
    \item The Lyapunov methods can be used to generate any arbitrary high-dimensional system, is it possible to add a benchmark with 128, 256, 512, or 1024 dimensional system? In reviewer's eyes $10$ may not be large system.
\end{enumerate}

}}

Please see Table \ref{tab:review-builtin-systems-20260314}.

\begin{table*}[ht]
\centering
\scriptsize
\setlength{\tabcolsep}{4pt}
\begin{tabular}{p{0.16\textwidth}p{0.07\textwidth}p{0.47\textwidth}p{0.16\textwidth}p{0.08\textwidth}}
\toprule
System & Dim. & Exact definition used in the paper & Reset distribution & Selected $dt$ \\
\midrule
Duffing & 2 & $x' = v$, $v' = x - x^3$ & $x$ uniform on $[-1.5, 1.5]$, $v$ uniform on $[-1, 1]$ & 0.01 \\
Lotka-Volterra & 2 & Predator-prey system with $\alpha = \beta = \gamma = \delta = 0.2$ & Both populations uniform on $[0.02, 3.0]$ & 0.0025 \\
Blended linear system & 2 & Softmax blend of three local linear systems centered at $(2,0)$, $(-2,2)$, and $(-2,-2)$ with local matrices having eigenvalues $-0.5 \pm 3i$, $(-0.2,-5.0)$, and $(-5.0,-2.0)$ respectively & Uniform on $[-4,4]^2$ & 0.05 \\
Multiwell gradient & 2 & Five canonical wells at $(-1,-1)$, $(1,-1)$, $(-1,1)$, $(1,1)$, and $(0,0)$; pure potential descent & Uniform on $[-2.5,2.5]^2$ & 0.02 \\
Multiwell rotational & 2 & Same wells as above; descent plus 0.6 times the 90-degree rotated gradient & Uniform on $[-2.5,2.5]^2$ & 0.02 \\
Multiwell energy & 2 & Same wells as above; descent plus radial forcing $1.2 \sin(r) x$ & Uniform on $[-2.5,2.5]^2$ & 0.02 \\
Multiwell strong transition & 2 & Same wells as above; descent plus 0.8 times the rotated gradient and radial forcing $\cos(2r) x$ & Uniform on $[-2.5,2.5]^2$ & 0.005 \\
Embedded multiwell gradient & 8 & Same 2D gradient dynamics in the first two coordinates, linear decay in the remaining six coordinates & Uniform on $[-2.5,2.5]^8$ & 0.005 \\
Embedded multiwell rotational & 8 & Same 2D rotational dynamics in the first two coordinates, linear decay in the remaining six coordinates & Uniform on $[-2.5,2.5]^8$ & 0.005 \\
Embedded multiwell energy & 8 & Same 2D energy dynamics in the first two coordinates, linear decay in the remaining six coordinates & Uniform on $[-2.5,2.5]^8$ & 0.02 \\
Embedded multiwell strong transition & 8 & Same 2D strong-transition dynamics in the first two coordinates, linear decay in the remaining six coordinates & Uniform on $[-2.5,2.5]^8$ & 0.005 \\
Kuramoto & 16 & Ring-coupled phase oscillators with periodic boundary conditions, coupling constant 4.0, and identical natural frequencies in the benchmark; basin label is the winding number & Each phase uniform on $[-\pi, \pi]$ & 0.0125 \\
Hopfield & 16 & Continuous Hopfield dynamics with 16 neurons, 4 stored random patterns in the benchmark, $\beta = 1.0$, $\tau = 1.0$, and zero bias; the stored patterns are random $\pm 1$ vectors and the recurrent weight matrix is $(P^\top P)/m$ with the diagonal reset to zero; basin label is maximal sign-overlap pattern index & Gaussian with standard deviation 1.0 & 0.0125 \\
Competitive Lotka-Volterra & 10 & Symmetric interaction matrix, interaction scale 0.70, uniform growth, positivity clipping, survival threshold $10^{-3}$; off-diagonal entries are sampled uniformly on $[0,0.70]$, symmetrized, and the diagonal is then set to 1.0; basin label is the survivor-support bitmask & Each coordinate uniform on $[0.05,1.5]$ & 0.01 \\
\bottomrule
\end{tabular}
\caption{\textbf{Built-in benchmark systems.}}
\label{tab:review-builtin-systems-20260314}
\end{table*}

\paragraph{Exact built-in equations used most often in the paper.}

The \textbf{blended linear system} is defined by
\begin{equation}
x' = \sum_{k=1}^{3} w_k(x)\, A_k (x - c_k),
\end{equation}
with
\begin{equation}
w_k(x) \propto \exp\!\left(-\frac{\|x-c_k\|^2}{2\sigma^2}\right),
\qquad
\sigma = 1.5.
\end{equation}

The \textbf{multiwell family} uses the scalar field
\begin{equation}
V(x) = -\sigma^2 \sum_j \exp\!\left(-\frac{\|x-c_j\|^2}{\sigma^2}\right),
\qquad
\sigma = 0.7,
\end{equation}
so that
\begin{equation}
\nabla V(x) = 2 \sum_j (x-c_j)\exp\!\left(-\frac{\|x-c_j\|^2}{\sigma^2}\right).
\end{equation}

The \textbf{four variants} are then defined by
\begin{equation}
x' = -\nabla V(x),
\end{equation}
\begin{equation}
x' = -\nabla V(x) + 0.6\, R_{90}\bigl(\nabla V(x)\bigr),
\end{equation}
\begin{equation}
x' = -\nabla V(x) + 1.2 \sin(\|x\|)\,x,
\end{equation}
and
\begin{equation}
x' = -\nabla V(x) + 0.8\, R_{90}\bigl(\nabla V(x)\bigr) + \cos(2\|x\|)\,x,
\end{equation}
where
\begin{equation}
R_{90}(a,b) = (-b,a).
\end{equation}

The \textbf{high-dimensional embedded variants} apply the corresponding two-dimensional core dynamics
to the first two coordinates and use
\begin{equation}
x_i' = -x_i
\end{equation}
for the remaining coordinates.

The \textbf{benchmark Kuramoto system} uses
\begin{equation}
\theta_i' = \omega_i + \frac{4.0}{N}
\left\{
\sin(\theta_{i+1}-\theta_i)
+
\sin(\theta_{i-1}-\theta_i)
\right\},
\end{equation}
with periodic boundary conditions and
\begin{equation}
\omega_i = 0
\qquad
\text{for all } i.
\end{equation}

\begin{table*}[ht]
\centering
\scriptsize
\setlength{\tabcolsep}{4pt}
\begin{tabular}{p{0.22\textwidth}p{0.08\textwidth}p{0.20\textwidth}p{0.20\textwidth}}
\toprule
Dysts system & Dim. & Native Dysts $dt$ & Selected benchmark $dt$ \\
\midrule
Dadras & 3 & 0.0006578296382730287 & 0.0006578296382730287 \\
Duffing & 3 & 0.0014901683613936711 & 0.0003725420903484178 \\
QiChen & 3 & $7.837106184364728 \times 10^{-5}$ & $7.837106184364728 \times 10^{-5}$ \\
Sakarya & 3 & 0.0009970461743625946 & 0.0009970461743625946 \\
SprottTorus & 3 & 0.0027218536422742544 & 0.0027218536422742544 \\
Chua & 3 & 0.0002847474579095888 & 0.0002847474579095888 \\
MultiChua & 3 & 0.00045179529921654196 & 0.00022589764960827098 \\
DequanLi & 3 & $1.6763993998996084 \times 10^{-5}$ & $4.190998499749021 \times 10^{-6}$ \\
LuChenCheng & 3 & 0.00018469678279714685 & 0.00018469678279714685 \\
SanUmSrisuchinwong & 3 & 0.0014933288881479711 & 0.0014933288881479711 \\
WangSun & 3 & 0.005392498749791912 & 0.001348124687447978 \\
ShimizuMorioka & 3 & 0.002408001333556058 & 0.0006020003333890145 \\
LorenzCoupled & 6 & 0.0003241940323387382 & $8.104850808468456 \times 10^{-5}$ \\
RikitakeDynamo & 3 & 0.0011796966161026034 & 0.00029492415402565086 \\
Hadley & 3 & 0.00029086847807974437 & 0.00029086847807974437 \\
\bottomrule
\end{tabular}
\caption{\textbf{Dysts benchmark systems.} All systems use the Dysts 0.96 default parameterization and the standardized reset rule described above.}
\label{tab:review-dysts-systems-20260314}
\end{table*}

\subsection{Cross-system forecasting family}

\paragraph{Aim and isolation logic.}
This family of experiments asks whether adding LISTA improves performance against a fair
sparse MLP anchor once training budget, latent dimension, front-end width and
depth, step-size table, seed set, and forecasting metric are matched. The
LISTA comparator is fixed before the final 29-system fair benchmark, so the
benchmark itself is not a per-system tuning exercise. 

\paragraph{Exact setup details.}
The family uses all 29 systems and 200,000 training steps per run. The sparse
MLP from \citet{fathi_course_2023}, the zero-sparsity MLP baseline, the standard LISTA comparator, and block-diagonal Koopman setups with the sparse MLP and LISTA encoder all
use matched seed coverage: seeds 0--9 on all 29 systems.
% The collected
% March 31 benchmark refresh below comes from
% \texttt{results/paper\_seed\_statistics\_20260331/forecasting\_aggregate\_stats.csv},
% \texttt{results/paper\_seed\_statistics\_20260331/forecasting\_builtin\_appendix\_seed\_stats.csv},
% and
% \texttt{results/paper\_seed\_statistics\_20260331/forecasting\_dysts\_appendix\_seed\_stats.csv}.
% Those files were recomputed from raw
% \texttt{evaluation\_results\_best.json} artifacts and verified against the
% original collector-row files in
% \texttt{docs/PAPER\_SEED\_STATISTICS\_20260331.md} before any IQM summary was
% formed. 
The collected
zero-sparsity benchmark extension keeps the exact same 200k recipe as the
sparse MLP anchor, except it sets $\lambda_{\mathrm{sparse}} = 0$. 

\paragraph{How to read the benchmark tables.}
The robust benchmark summary separates two notions of strength. The
cross-system IQM asks which model is better after treating each system-level
seed mean as one unit and trimming the heaviest tails (outliers). The win count asks on
how many of the 29 matched systems one model has the lower system-level seed
IQM. 
Selected bootstrap
intervals are reported in Table \ref{tab:review-cross-system-iqm-ci-20260314}.
Those intervals overlap broadly on the headline benchmark, so small IQM ordering changes should be read cautiously.

\paragraph{System-by-system benchmark read.}
The split-horizon story is not a one-system artifact. Tables
\ref{tab:review-cross-system-builtin-20260314} and
\ref{tab:review-cross-system-dysts-20260314} report the per-system seed-median
best-periodic errors for the two paper-facing benchmark roots at $H1000$ and
$H3000$. At $H1000$, standard LISTA is lower on 18 of 29 systems, but the
sparse MLP still keeps the better cross-system median because its losses are
more concentrated and usually smaller on the systems where it wins. By
$H3000$, standard LISTA still wins 18 of 29 systems, but the remaining sparse
MLP wins include the safer rows: for example, Hopfield, the embedded
strong-transition multiwell, and multiwell-gradient HD all favor the sparse MLP
by a large margin.





\paragraph{System-by-system benchmark read.}
The split-horizon story is not a one-system artifact. Tables
\ref{tab:review-cross-system-dysts-20260314} and
\ref{tab:review-cross-system-builtin-20260314} report the per-system seed-IQM
best-periodic errors for the three paper-facing benchmark roots at $H1000$ and
$H3000$; the corresponding bootstrap 95\% confidence intervals are in the
appendix CSVs cited in the captions. 

The Dysts subset is now strongly
Dense-LISTA-positive: dense LISTA wins 12/15 Dysts systems at $H1000$ and
13/15 at $H3000$, with DequanLi and SprottTorus as the durable sparse-MLP
exceptions and Dysts Duffing as a zero-MLP exception only at $H1000$. 

The
built-in systems remain mixed: the zero-sparsity MLP is still strongest on
blended, benchmark Hopfield, the canonical Kuramoto replacement, the canonical
competitive-LV follow-up row, and several embedded multiwell settings, while
dense LISTA remains clearly better on the standard multiwell systems and the
strong-transition row. This is the first sign that the
sparsity penalty is not acting like a generic regularizer: it helps some
systems by sharpening separable support structure, but it can also obstruct
forecasting on systems where useful latent structure is more continuous or more
distributed than a sparse-code story would suggest. 
% To avoid mixing two
% different CLV follow-up conventions, the built-in table now uses the fixed
% 8-basin competitive Lotka--Volterra $dt=0.005$ row as its canonical CLV entry.
% The built-in table now also uses the completed canonical Kuramoto replacement:
% $N=16$, identical frequencies, $dt=0.01$, 200k, from
% \texttt{results/kuramoto\_dt0p01\_200k\_canonical\_20260323}, so the built-in
% table and the dedicated Kuramoto discussion now refer to the same setting.

\begin{table*}[ht]
\centering
\scriptsize
\setlength{\tabcolsep}{4pt}
\begin{tabular}{p{0.1\textwidth}rrrrrr}
\toprule
Dysts system & Sparse $H1000$ IQM & Zero $H1000$ IQM & LISTA $H1000$ IQM & Sparse $H3000$ IQM & Zero $H3000$ IQM & LISTA $H3000$ IQM \\
\midrule
Dadras & 0.0111 & 0.0102 & \textbf{0.0096} & 0.4812 & 0.5590 & \textbf{0.3864} \\
Dysts Duffing & 0.8542 & \textbf{0.8488} & 0.8522 & 2.1228 & 1.9001 & \textbf{1.8717} \\
QiChen & 0.0042 & 0.0052 & \textbf{0.0029} & 0.0770 & 0.1753 & \textbf{0.0492} \\
Sakarya & 0.0128 & 0.0117 & \textbf{0.0091} & 0.2266 & 0.1986 & \textbf{0.1700} \\
SprottTorus & 0.0085 & \textbf{0.0069} & 0.0089 & \textbf{0.0914} & 0.1484 & 0.1102 \\
Chua & 0.0018 & 0.0017 & \textbf{8.786e-04} & 0.0173 & 0.0193 & \textbf{0.0071} \\
MultiChua & 0.0034 & 0.0048 & \textbf{0.0022} & 0.0386 & 0.0556 & \textbf{0.0331} \\
DequanLi & \textbf{0.0048} & 0.0110 & 0.0050 & \textbf{0.0404} & 0.1388 & 0.0472 \\
LuChenCheng & 0.0018 & 0.0024 & \textbf{8.778e-04} & 0.0176 & 0.0214 & \textbf{0.0062} \\
SanUmSrisuchinwong & 0.0056 & 0.0040 & \textbf{0.0038} & 0.0774 & 0.0399 & \textbf{0.0354} \\
WangSun & 0.0281 & 0.0557 & \textbf{0.0158} & 0.6715 & 1.2392 & \textbf{0.5567} \\
Shimizu--Morioka & 0.0044 & 0.0071 & \textbf{0.0016} & 0.0706 & 0.0661 & \textbf{0.0214} \\
LorenzCoupled & 0.0272 & 0.0353 & \textbf{0.0257} & 0.9806 & 1.2007 & \textbf{0.7922} \\
Rikitake dynamo & 0.0038 & 0.0075 & \textbf{0.0028} & 0.0826 & 0.1091 & \textbf{0.0402} \\
Hadley & 0.0040 & 0.0035 & \textbf{0.0032} & 0.0321 & 0.0330 & \textbf{0.0230} \\
\midrule
Wins & 1/15 & 2/15 & \textbf{12/15} & 2/15 & 0/15 & \textbf{13/15} \\
\bottomrule
\end{tabular}
\caption{\textbf{Dysts system-by-system benchmark comparison.} Entries are system-level seed IQMs of the best-periodic forecasting error, copied from \texttt{results/paper\_seed\_statistics\_20260331/forecasting\_dysts\_appendix\_seed\_stats.csv}. Bootstrap 95\% confidence intervals for each cell are in the same CSV. Bold indicates the lower seed-IQM value at that horizon among the sparse anchor, the matched zero-sparsity MLP, and the dense LISTA comparator. The last row reports within-Dysts system-win counts at each horizon.}
\label{tab:review-cross-system-dysts-20260314}
\end{table*}

\begin{table*}[ht]
\centering
\scriptsize
\setlength{\tabcolsep}{4pt}
\begin{tabular}{p{0.26\textwidth}rrrrrr}
\toprule
Built-in system & Sparse $H1000$ IQM & Zero $H1000$ IQM & LISTA $H1000$ IQM & Sparse $H3000$ IQM & Zero $H3000$ IQM & LISTA $H3000$ IQM \\
\midrule
Lotka--Volterra & \textbf{0.0014} & 0.0035 & 0.0019 & \textbf{0.0177} & 0.0243 & 0.1074 \\
Duffing & \textbf{0.0285} & 0.0586 & 0.0370 & \textbf{0.2810} & 0.4169 & 0.3670 \\
Blended linear & 0.1168 & \textbf{0.0492} & 0.0706 & 0.1265 & \textbf{0.0498} & 0.0729 \\
Multiwell gradient & 0.0548 & 0.0738 & \textbf{0.0316} & 4.2817 & 0.2000 & \textbf{0.0639} \\
Multiwell rotational & 0.0669 & 0.0811 & \textbf{0.0464} & 0.1412 & 0.1877 & \textbf{0.0847} \\
Multiwell energy & \textbf{0.5019} & 0.5141 & 0.5198 & 0.8786 & 1.0227 & \textbf{0.7923} \\
Multiwell strong transition & 1.068e+28 & 2.997e+20 & \textbf{0.3962} & 0.7719 & 1.826e+05 & \textbf{0.2024} \\
Embedded multiwell gradient & 1.4141 & \textbf{1.0659} & 1.3654 & 7992.539 & \textbf{1.4823} & 11.7782 \\
Embedded multiwell rotational & 2.1160 & \textbf{1.2052} & 4.4237 & 1.637e+04 & \textbf{1.5664} & 3.000e+05 \\
Embedded multiwell energy & \textbf{1.9597} & 2.4141 & 5.1342 & \textbf{2.4634} & 3.1888 & 6.2702 \\
Embedded multiwell strong transition & 1.111e+31 & \textbf{1.3350} & 1.137e+04 & 1.011e+07 & \textbf{1.3544} & 2.033e+10 \\
Kuramoto & 682.989 & \textbf{39.9860} & 82.5883 & 1.799e+11 & \textbf{4.377e+06} & 4.770e+11 \\
Hopfield & 233.598 & \textbf{57.4201} & 7.738e+16 & 8.961e+05 & \textbf{73.8126} & 3.784e+26 \\
Competitive LV & 0.5259 & \textbf{0.4027} & 0.5696 & 0.6352 & \textbf{0.4847} & 0.6791 \\
\midrule
Wins & 4/14 & \textbf{7/14} & 3/14 & 3/14 & \textbf{7/14} & 4/14 \\
\bottomrule
\end{tabular}
\caption{\textbf{Built-in system-by-system benchmark comparison.} Entries are system-level seed IQMs of the best-periodic forecasting error.
% , copied from \texttt{results/paper\_seed\_statistics\_20260331/forecasting\_builtin\_appendix\_seed\_stats.csv}
Bootstrap 95\% confidence intervals for each cell are in the same CSV. Bold indicates the lower seed-IQM value at that horizon among the sparse anchor, the matched zero-sparsity MLP, and the dense LISTA comparator. The Kuramoto row uses the completed canonical $N=16$, identical-frequency, $dt=0.01$ follow-up, and the competitive Lotka--Volterra row uses the fixed 8-basin $dt=0.005$ canonical follow-up values. The last row reports within-built-in system-win counts at each horizon.}
\label{tab:review-cross-system-builtin-20260314}
\end{table*}

\begin{table*}[ht]
\centering
\scriptsize
\setlength{\tabcolsep}{5pt}
\begin{tabular}{lrrrrrrrr}
\toprule
Horizon & Sparse MLP IQM & Zero MLP IQM & LISTA IQM & Zero wins & LISTA wins & Good sparse & Good zero & Good LISTA \\
\midrule
$H100$  & \textbf{0.0034} & 0.0038 & 0.0037 & 13/29 & 10/29 & 27/29 & 27/29 & 28/29 \\
$H500$  & 0.0585 & 0.0610 & \textbf{0.0548} & 11/29 & 17/29 & 25/29 & 26/29 & 25/29 \\
$H1000$ & 0.2434 & 0.1480 & \textbf{0.0910} & 14/29 & 20/29 & 25/29 & 26/29 & 26/29 \\
$H1500$ & 0.3276 & 0.2796 & \textbf{0.1765} & 10/29 & 22/29 & 26/29 & 26/29 & 25/29 \\
$H2000$ & \textbf{0.3108} & 0.3865 & 0.3574 & 10/29 & 21/29 & 24/29 & 26/29 & 25/29 \\
$H2500$ & 0.4090 & 0.4927 & \textbf{0.3105} & 11/29 & 20/29 & 24/29 & 26/29 & 25/29 \\
$H3000$ & \textbf{0.5068} & 0.5902 & 0.5725 & 11/29 & 20/29 & 24/29 & 26/29 & 24/29 \\
\bottomrule
\end{tabular}
\caption{\textbf{Main 29-system benchmark result.} The displayed values are cross-system IQMs over the 29 system-level seed means from \texttt{results/paper\_seed\_statistics\_20260331/forecasting\_aggregate\_stats.csv}. `Zero wins' and `LISTA wins' count systems where that model beats the sparse MLP anchor on the system-level seed IQM. `Good' counts use system-level seed IQM $< 10$. Bootstrap 95\% confidence intervals are in the same CSV.}
\label{tab:review-cross-system-main-20260314}
\end{table*}


\paragraph{Fixed-cadence diagnostic.}
The main benchmark uses horizon-wise best periodic re-encoding, so it is an
evaluation-time oracle over a fixed candidate set. A useful deployment-like
diagnostic is to force one global cadence for every system and horizon. Table
\ref{tab:review-cross-system-fixed-cadence-20260314} shows that under one fixed
\texttt{periodic\_100} policy, standard LISTA still wins more systems at the long
horizons, but the clean late-horizon median reversal is weaker: the sparse MLP
is better by cross-system median at $H2000$ and $H2500$, and it keeps better
$H3000$ coverage.

\begin{table*}[ht]
\centering
\scriptsize
\setlength{\tabcolsep}{5pt}
\begin{tabular}{lrrrrrrrr}
\toprule
Horizon & Sparse MLP median & Zero MLP median & LISTA median & Zero wins & LISTA wins & Good sparse & Good zero & Good LISTA \\
\midrule
$H1500$ & 0.0749 & 0.1726 & \textbf{0.0535} & 8/29 & \textbf{18/29} & 19/29 & 19/29 & \textbf{20/29} \\
$H2000$ & \textbf{0.1829} & 0.4713 & 0.2787 & 10/29 & \textbf{17/29} & \textbf{20/29} & 19/29 & \textbf{20/29} \\
$H2500$ & \textbf{0.3417} & 0.7047 & 0.4447 & 8/29 & \textbf{17/29} & \textbf{20/29} & 19/29 & 19/29 \\
$H3000$ & \textbf{0.5930} & 0.8987 & 0.6182 & 10/29 & \textbf{19/29} & \textbf{20/29} & 19/29 & 17/29 \\
\bottomrule
\end{tabular}
\caption{\textbf{Fixed-cadence diagnostic using one global re-encoding period of 100 steps.} This is not the official benchmark rule; it is a supporting diagnostic on the same saved checkpoints.}
\label{tab:review-cross-system-fixed-cadence-20260314}
\end{table*}




\begin{table*}[ht]
\centering
\tiny
\setlength{\tabcolsep}{4pt}
\begin{tabular}{p{0.24\textwidth}ccc}
\toprule
Dysts system & Sparse $H1000$ 95\% CI & Zero $H1000$ 95\% CI & LISTA $H1000$ 95\% CI \\
\midrule
Dadras & [0.0063, 0.0229] & [0.0065, 0.0214] & [0.0067, 0.0163] \\
Dysts Duffing & [0.8102, 0.9453] & [0.7887, 0.9525] & [0.7788, 0.9254] \\
QiChen & [0.0032, 0.0049] & [0.0037, 0.0086] & [0.0025, 0.0037] \\
Sakarya & [0.0086, 0.0319] & [0.0065, 0.0487] & [0.0054, 0.0317] \\
SprottTorus & [0.0059, 0.0102] & [0.0041, 0.0130] & [0.0075, 0.0125] \\
Chua & [0.0012, 0.0026] & [7.236e-04, 0.0057] & [6.136e-04, 0.0011] \\
MultiChua & [0.0027, 0.0050] & [0.0031, 0.0077] & [0.0018, 0.0033] \\
DequanLi & [0.0038, 0.0063] & [0.0082, 0.0206] & [0.0039, 0.0060] \\
LuChenCheng & [0.0013, 0.0021] & [0.0011, 0.0037] & [6.145e-04, 0.0011] \\
SanUmSrisuchinwong & [0.0038, 0.0081] & [0.0028, 0.0071] & [0.0028, 0.0051] \\
WangSun & [0.0132, 0.0523] & [0.0239, 0.0885] & [0.0106, 0.0275] \\
Shimizu--Morioka & [0.0036, 0.0065] & [0.0033, 0.0128] & [0.0012, 0.0038] \\
LorenzCoupled & [0.0214, 0.0398] & [0.0250, 0.0541] & [0.0235, 0.0293] \\
Rikitake dynamo & [0.0027, 0.0092] & [0.0056, 0.0104] & [0.0017, 0.0120] \\
Hadley & [0.0034, 0.0043] & [0.0020, 0.0066] & [0.0027, 0.0040] \\
\bottomrule
\end{tabular}
\caption{\textbf{Dysts system-by-system bootstrap intervals at $H1000$.} Entries are bootstrap 95\% confidence intervals for the seed-IQM values reported in Table \ref{tab:review-cross-system-dysts-20260314}, copied from \texttt{results/paper\_seed\_statistics\_20260331/forecasting\_dysts\_appendix\_seed\_stats.csv}.}
\label{tab:review-cross-system-dysts-ci-h1000-20260314}
\end{table*}

\begin{table*}[ht]
\centering
\tiny
\setlength{\tabcolsep}{4pt}
\begin{tabular}{p{0.24\textwidth}ccc}
\toprule
Dysts system & Sparse $H3000$ 95\% CI & Zero $H3000$ 95\% CI & LISTA $H3000$ 95\% CI \\
\midrule
Dadras & [0.2877, 0.7787] & [0.3424, 0.7705] & [0.1971, 0.5787] \\
Dysts Duffing & [1.8834, 2.6231] & [1.8033, 2.0922] & [1.8140, 2.0060] \\
QiChen & [0.0379, 0.2087] & [0.0870, 0.3177] & [0.0373, 0.0984] \\
Sakarya & [0.1522, 0.3608] & [0.1187, 0.3724] & [0.0982, 0.2856] \\
SprottTorus & [0.0563, 0.1453] & [0.0883, 0.2756] & [0.0833, 0.1539] \\
Chua & [0.0102, 0.0291] & [0.0083, 0.0839] & [0.0050, 0.0138] \\
MultiChua & [0.0268, 0.0658] & [0.0319, 0.1242] & [0.0260, 0.0504] \\
DequanLi & [0.0309, 0.0638] & [0.0761, 0.2207] & [0.0391, 0.0515] \\
LuChenCheng & [0.0117, 0.0363] & [0.0120, 0.0376] & [0.0041, 0.0137] \\
SanUmSrisuchinwong & [0.0410, 0.1670] & [0.0244, 0.0976] & [0.0222, 0.0577] \\
WangSun & [0.3526, 1.0998] & [0.8907, 1.3865] & [0.2550, 0.8683] \\
Shimizu--Morioka & [0.0423, 0.1703] & [0.0346, 0.1224] & [0.0158, 0.0387] \\
LorenzCoupled & [0.7446, 1.3090] & [0.8025, 1.5789] & [0.6803, 1.0112] \\
Rikitake dynamo & [0.0366, 0.1606] & [0.0744, 0.1618] & [0.0145, 0.0866] \\
Hadley & [0.0272, 0.0377] & [0.0150, 0.0617] & [0.0175, 0.0320] \\
\bottomrule
\end{tabular}
\caption{\textbf{Dysts system-by-system bootstrap intervals at $H3000$.} Entries are bootstrap 95\% confidence intervals for the seed-IQM values reported in Table \ref{tab:review-cross-system-dysts-20260314}, copied from \texttt{results/paper\_seed\_statistics\_20260331/forecasting\_dysts\_appendix\_seed\_stats.csv}.}
\label{tab:review-cross-system-dysts-ci-h3000-20260314}
\end{table*}

\begin{table*}[ht]
\centering
\tiny
\setlength{\tabcolsep}{4pt}
\begin{tabular}{p{0.24\textwidth}ccc}
\toprule
Built-in system & Sparse $H1000$ 95\% CI & Zero $H1000$ 95\% CI & LISTA $H1000$ 95\% CI \\
\midrule
Lotka--Volterra & [8.638e-04, 0.0021] & [0.0014, 0.0076] & [8.110e-04, 0.0038] \\
Duffing & [0.0198, 0.0366] & [0.0447, 0.0710] & [0.0264, 0.0574] \\
Blended linear & [0.0213, 0.3610] & [0.0079, 0.2338] & [0.0130, 0.1969] \\
Multiwell gradient & [0.0288, 5.5118] & [0.0499, 0.3712] & [0.0192, 0.0457] \\
Multiwell rotational & [0.0418, 0.0880] & [0.0655, 0.1145] & [0.0334, 0.0635] \\
Multiwell energy & [0.4000, 0.6422] & [0.4156, 0.6306] & [0.3087, 0.7686] \\
Multiwell strong transition & [5.523e+04, 7.221e+30] & [1.762e+18, 6.231e+26] & [0.0764, 0.8383] \\
Embedded multiwell gradient & [0.9977, 19.0433] & [0.9074, 1.1747] & [1.0842, 8.5994] \\
Embedded multiwell rotational & [1.3007, 3.9895] & [1.0772, 1.3158] & [1.2825, 12.4121] \\
Embedded multiwell energy & [1.7155, 2.5357] & [1.7536, 2.9310] & [4.5683, 5.8636] \\
Embedded multiwell strong transition & [2.353e+26, 1.653e+33] & [1.0898, 1.7359] & [101.585, 5.599e+18] \\
Kuramoto & [534.476, 824.822] & [28.1987, 65.6480] & [54.5185, 1420.344] \\
Hopfield & [128.571, 2.672e+04] & [44.8407, 81.9247] & [4.224e+09, 1.508e+24] \\
Competitive LV & [0.4564, 0.5997] & [0.3601, 0.4781] & [0.5184, 0.6164] \\
\bottomrule
\end{tabular}
\caption{\textbf{Built-in system-by-system bootstrap intervals at $H1000$.} Entries are bootstrap 95\% confidence intervals for the seed-IQM values reported in Table \ref{tab:review-cross-system-builtin-20260314}, copied from \texttt{results/paper\_seed\_statistics\_20260331/forecasting\_builtin\_appendix\_seed\_stats.csv}.}
\label{tab:review-cross-system-builtin-ci-h1000-20260314}
\end{table*}

\begin{table*}[ht]
\centering
\tiny
\setlength{\tabcolsep}{4pt}
\begin{tabular}{p{0.24\textwidth}ccc}
\toprule
Built-in system & Sparse $H3000$ 95\% CI & Zero $H3000$ 95\% CI & LISTA $H3000$ 95\% CI \\
\midrule
Lotka--Volterra & [0.0115, 0.0327] & [0.0142, 0.0647] & [0.0503, 3.1457] \\
Duffing & [0.2188, 0.4421] & [0.3376, 0.4974] & [0.2741, 0.5966] \\
Blended linear & [0.0159, 0.4043] & [0.0040, 0.2638] & [0.0052, 0.2073] \\
Multiwell gradient & [0.1653, 2.794e+18] & [0.0921, 3.279e+04] & [0.0344, 0.0983] \\
Multiwell rotational & [0.0993, 0.2008] & [0.1467, 0.2457] & [0.0571, 3.783e+10] \\
Multiwell energy & [0.6283, 1.1546] & [0.7741, 1.4442] & [0.5581, 1.1167] \\
Multiwell strong transition & [0.1608, 1.4284] & [0.1666, 3.642e+31] & [0.0474, 0.6087] \\
Embedded multiwell gradient & [3.1326, 9.440e+11] & [1.3077, 1.5758] & [1.9563, 50.6341] \\
Embedded multiwell rotational & [1.8430, 1.310e+07] & [1.4280, 1.6319] & [15.6501, 4.116e+08] \\
Embedded multiwell energy & [2.0969, 3.3914] & [2.5743, 3.7426] & [5.6247, 7.9544] \\
Embedded multiwell strong transition & [1.4654, 1.151e+24] & [1.1040, 1.7831] & [37.2120, 3.270e+20] \\
Kuramoto & [6.548e+10, 4.131e+11] & [6.980e+05, 1.036e+08] & [1.575e+08, 3.637e+19] \\
Hopfield & [170.274, 3.113e+11] & [61.3652, 100.513] & [1.254e+19, 9.247e+32] \\
Competitive LV & [0.5593, 0.7101] & [0.4405, 0.5771] & [0.6282, 0.7186] \\
\bottomrule
\end{tabular}
\caption{\textbf{Built-in system-by-system bootstrap intervals at $H3000$.} Entries are bootstrap 95\% confidence intervals for the seed-IQM values reported in Table \ref{tab:review-cross-system-builtin-20260314}, copied from \texttt{results/paper\_seed\_statistics\_20260331/forecasting\_builtin\_appendix\_seed\_stats.csv}.}
\label{tab:review-cross-system-builtin-ci-h3000-20260314}
\end{table*}

\subsection{Hard-system forecasting family}

\paragraph{Systems and why they are hard.}
This family focuses on three benchmark settings that are harder than the
typical systems in the main 29-system table because they combine higher
dimensional state spaces with richer multi-basin structure. Kuramoto is a
phase-oscillator system with 16--64 coupled phases and winding-number basin
structure. Hopfield uses 16 or 64 continuous neurons with multiple attractors
induced by stored patterns. The repaired competitive Lotka--Volterra controls
combine a corrected 4-basin setting and a fixed 8-basin setting, where
distinct survivor-support patterns define multiple long-run outcomes. These
systems are therefore useful for asking whether the same models that work well
on the broader benchmark still forecast accurately when the state dimension is
higher, the number of basins is larger, and the basin geometry is less local.

\paragraph{Aim and isolation logic.}
The goal of this family is to compare model behavior on these harder systems
under controlled interventions. The main interventions are smaller integration
step size, encoder family, and latent-transition structure. Within each
targeted setting, all compared models use the same system, same step size, same
training budget, same latent dimension, and same evaluation protocol. Where
needed, a block-diagonal MLP control isolates the effect of block structure in
the latent transition from the effect of the encoder family.

Hyperparameter tuning was done with a matched tuning budget across model
families. The first stage used 6 coarse configurations per family per setting,
100000 training steps, and seeds 0 and 1. The second stage re-ran one selected
winner per family per setting for 200000 steps with seeds 0, 1, and 2.
The March 21 zero-sparsity MLP extension is now collected on every
paper-facing hard-system setting in this section with the same system, step
size, horizon table, and seed coverage as the sparse-MLP rows, but with
$\lambda_{\mathrm{sparse}} = 0$. The March 22 recovery also refreshed the
Kuramoto $N=16$ identical collector so the zero-sparsity rows key correctly as
Kuramoto rather than an intermediate dimension label. The hard-system tables
below are copied from
\texttt{results/paper\_seed\_statistics\_20260331/forecasting\_hard\_system\_seed\_stats.csv},
which is the March 31 raw-artifact reanalysis verified against the original
hard-system collector rows in
\texttt{docs/PAPER\_SEED\_STATISTICS\_20260331.md}.

\paragraph{Main empirical read.}
The cleanest summary is still system-dependent, but the IQM tables narrow the
positive claims further. On Kuramoto, smaller $dt$ plus block structure remains
the strongest long-horizon story at $N=16/N24/N32$, while the zero-sparsity
MLP becomes the strongest row at $N=64$. On Hopfield, the zero-sparsity MLP is
best at $H100$ and $H1000$--$H3000$ on the targeted $N=16$ setting, but the
sparse MLP remains slightly lower at $H500$; in the quarter-step $N=64$ probe,
the sparse anchor is best only at $H100/H500$ and the zero-sparsity MLP takes
over from $H1000$ onward. On competitive Lotka--Volterra, removing the
sparsity penalty now wins the corrected 4-basin table at every horizon and
also wins the fixed 8-basin $dt=0.005$ table at every horizon under IQM,
whereas the smaller-step $dt=0.0025$ follow-up still favors the sparse anchor
through $H2500$ before the zero-sparsity MLP edges it at $H3000$. The purpose
of this section is therefore not to claim a uniform hard-system advantage for
one method, but to document that explicit sparsity can help some systems while
clearly hurting others once forecasting difficulty increases. This is also why
the paper should avoid the broader sentence ``sparsity is good for
long-horizon forecasting'': the hard-system family contains several direct
counterexamples, especially on Hopfield and competitive Lotka--Volterra, where
removing the penalty improves long-horizon forecasting substantially.

\paragraph{Exact parity grids.}
For Hopfield, the sparse MLP grid uses step size in
\{0.00625, 0.003125, 0.0015625\} and sparsity weight in \{0.0005, 0.0025\}. The
standard LISTA grid uses the same three step sizes and two recipes: the promoted benchmark recipe and a legacy parity recipe with sparsity weight 0.006,
main learning rate $10^{-4}$, and latent-transition learning rate $10^{-5}$.
The block-diagonal LISTA grid uses the same three step sizes and two recipes: a
targeted hard-system recipe with sparsity weight 0.001, one LISTA iteration,
$\alpha = 0.15$, and block size 16, and a transferred parity-style recipe with
sparsity weight 0.003, one LISTA iteration, $\alpha = 0.15$, and block size 16.
The block-diagonal MLP fairness-control grid mirrors the block-diagonal LISTA
non-encoder settings exactly. For competitive Lotka--Volterra, the corrected
4-basin rebuild reports the full 15-seed 200k collector at $dt=0.01$, and the
fixed 8-basin follow-up reports the $dt=0.005$ and $dt=0.0025$ reruns with the
same matched control logic. Other competitive Lotka--Volterra settings were
checked, but they are omitted here because they do not change the paper-facing
ordering.

Both repaired competitive Lotka--Volterra families use the same
$N$-species ODE
$x_i' = x_i (r_i - \sum_j A_{ij} x_j)$
with $r_i = 1$, $A_{ii} = 1$, symmetric off-diagonal interactions, RK4
integration, positivity clipping after each step, reset
$x_i(0) \sim \mathrm{Uniform}[0.05, 1.5]$, and survivor-support bitmask labels
at threshold $10^{-3}$. The corrected 4-basin rebuild uses $N=10$,
interaction scale $s=0.70$, and a training-seed-dependent system realization, so the interaction
matrix varies with the training seed; `4-basin' is empirical shorthand for the
observed major survivor-support count of that benchmark configuration. The
fixed 8-basin follow-up uses $N=15$, interaction scale $s=0.83$, and one fixed
system realization shared across seeds; `8-basin' is likewise the observed major
survivor-support count of that fixed system, not a quantity assumed known at
training time.


\begin{table*}[ht]
\centering
\scriptsize
\setlength{\tabcolsep}{3.5pt}
\begin{tabular}{p{0.19\textwidth}p{0.16\textwidth}rrrrrrr}
\toprule
Setting & Model & H100 & H500 & H1000 & H1500 & H2000 & H2500 & H3000 \\
\midrule
\multicolumn{9}{l}{\textit{Corrected 4-basin competitive Lotka--Volterra, full 15-seed rebuild, $dt=0.01$, 200k}} \\
& Sparse MLP & 0.0111 & 0.1110 & 0.1520 & 0.1820 & 0.2082 & 0.2274 & 0.2441 \\
& Zero MLP & \textbf{0.0085} & \textbf{0.0951} & \textbf{0.1285} & \textbf{0.1506} & \textbf{0.1710} & \textbf{0.1892} & \textbf{0.2058} \\
& Block-diag.\ MLP & 0.0128 & 0.2243 & 0.2891 & 0.3333 & 0.3713 & 0.4064 & 0.4364 \\
& Dense LISTA & 0.0122 & 0.1109 & 0.1477 & 0.1674 & 0.1864 & 0.2037 & 0.2199 \\
& Block-diag.\ LISTA & 0.0567 & 7.318e+30 & 0.2436 & 0.2779 & 0.3169 & 0.3534 & 0.3851 \\
\midrule
\multicolumn{9}{l}{\textit{Fixed 8-basin competitive Lotka--Volterra, $dt=0.005$, 200k}} \\
& Sparse MLP & 0.0779 & 0.4050 & 0.5259 & 0.5740 & 0.6014 & 0.6204 & 0.6352 \\
& Zero MLP & \textbf{0.0484} & \textbf{0.3107} & \textbf{0.4027} & \textbf{0.4382} & \textbf{0.4585} & \textbf{0.4730} & \textbf{0.4847} \\
& Block-diag.\ MLP & 0.1343 & 0.5561 & 0.5484 & 0.5910 & 0.6263 & 0.6501 & 0.6690 \\
& Dense LISTA & 0.0567 & 0.4353 & 0.5696 & 0.6203 & 0.6476 & 0.6656 & 0.6791 \\
& Block-diag.\ LISTA & 0.0776 & 1.266e+15 & 9.537e+32 & 9.537e+32 & 9.537e+32 & 9.537e+32 & 9.537e+32 \\
\midrule
\multicolumn{9}{l}{\textit{Fixed 8-basin competitive Lotka--Volterra, $dt=0.0025$, 200k}} \\
& Sparse MLP & \textbf{0.0026} & \textbf{0.1197} & \textbf{0.2518} & \textbf{0.3193} & \textbf{0.3582} & \textbf{0.3841} & 0.4029 \\
& Zero MLP & 0.0105 & 0.1748 & 0.2876 & 0.3391 & 0.3676 & 0.3857 & \textbf{0.3983} \\
& Block-diag.\ MLP & 0.0116 & 0.2892 & 0.4899 & 0.5780 & 0.6269 & 0.6590 & 0.6820 \\
& Dense LISTA & 0.0107 & 0.2742 & 0.5233 & 0.6306 & 0.6516 & 0.6602 & 0.6831 \\
& Block-diag.\ LISTA & 0.0199 & 0.3223 & 0.5317 & 0.6222 & 0.6645 & 0.6913 & 0.7101 \\
\bottomrule
\end{tabular}
\caption{\textbf{Repaired competitive Lotka--Volterra fairness and follow-up results.} Entries are seed-IQM best-periodic errors from \texttt{results/paper\_seed\_statistics\_20260331/forecasting\_hard\_system\_seed\_stats.csv}, with the best value in each system-setting and horizon shown in bold. Bootstrap 95\% confidence intervals are in the same CSV.}
\label{tab:review-hard-clv-combined-20260314}
\end{table*}

\begin{table*}[ht]
\centering
\scriptsize
\setlength{\tabcolsep}{3.5pt}
\begin{tabular}{p{0.26\textwidth}p{0.15\textwidth}rrrrrrr}
\toprule
Setting & Model & H100 & H500 & H1000 & H1500 & H2000 & H2500 & H3000 \\
\midrule
\multicolumn{9}{l}{\textit{Kuramoto, $N=16$, identical frequencies, $dt=0.00625$, 200k}} \\
& Sparse MLP & \textbf{0.0401} & \textbf{1.5850} & 30.6451 & 552.540 & 1.216e+04 & 3.155e+05 & 9.128e+06 \\
& Zero MLP & 0.1320 & 2.7100 & 9.1145 & 26.8238 & 104.769 & 533.433 & 3212.600 \\
& Block-diag.\ MLP & 0.1280 & 2.4325 & \textbf{6.2799} & \textbf{9.6448} & \textbf{12.4115} & \textbf{14.7781} & \textbf{17.0724} \\
& Dense LISTA & 0.2221 & 4.4524 & 16.8187 & 74.2240 & 686.105 & 5605.660 & 5.462e+04 \\
& Block-diag.\ LISTA & 0.1537 & 2.6974 & 7.0325 & 10.9739 & 14.5291 & 17.9909 & 21.7155 \\
\midrule
\multicolumn{9}{l}{\textit{Kuramoto, $N=16$, uniform frequency spread, $dt=0.00625$, 200k}} \\
& Sparse MLP & \textbf{0.0339} & \textbf{1.7644} & 41.0754 & 854.873 & 2.084e+04 & 5.458e+05 & 1.487e+07 \\
& Zero MLP & 0.1362 & 2.7220 & 8.8096 & 23.5269 & \textbf{79.1980} & \textbf{353.974} & 1923.638 \\
& Block-diag.\ MLP & 0.1350 & 2.5708 & \textbf{8.0214} & \textbf{22.1277} & 85.6596 & 359.818 & \textbf{1433.056} \\
& Dense LISTA & 0.2128 & 4.2819 & 17.8496 & 67.1720 & 500.736 & 9195.673 & 1.680e+05 \\
& Block-diag.\ LISTA & 0.1648 & 2.9840 & 10.3918 & 42.5226 & 256.748 & 1626.861 & 1.007e+04 \\
\midrule
\multicolumn{9}{l}{\textit{Kuramoto, $N=32$, identical frequencies, $dt=0.00625$, 200k}} \\
& Sparse MLP & \textbf{0.1159} & 2.0689 & 7.2005 & 14.6645 & 27.7706 & 52.1447 & 113.327 \\
& Block-diag.\ MLP & 0.1381 & \textbf{1.6532} & \textbf{5.2075} & \textbf{9.3354} & \textbf{13.3994} & \textbf{17.0582} & \textbf{20.4151} \\
& Dense LISTA & 52.4291 & 76.7164 & 91.7743 & 100.962 & 108.690 & 114.777 & 121.072 \\
& Block-diag.\ LISTA & 0.2112 & 1.9757 & 6.1287 & 11.2966 & 16.7969 & 22.7895 & 29.9114 \\
\bottomrule
\end{tabular}
\caption{\textbf{Kuramoto forecasting results.} All entries are seed-IQM best-periodic errors from \texttt{results/paper\_seed\_statistics\_20260331/forecasting\_hard\_system\_seed\_stats.csv}. The $N=16$ identical-frequency table is the main targeted comparison; the repaired block-diagonal MLP control confirms that the long-horizon gain is not only an encoder-family effect. In the uniform-spread row, the repaired block-diagonal MLP remains best at $H1000/H1500/H3000$, while the zero-sparsity MLP is slightly lower at $H2000/H2500$. In the $N=32$ row, the sparse-MLP and block-diagonal-LISTA entries remain the dedicated March 9 moderate-dimension confirmation, while the recovered block-diagonal-MLP and dense-LISTA entries come from the matched four-model dimension sweep at the same setting. Bootstrap 95\% confidence intervals are in the same CSV.}
\label{tab:review-hard-kuramoto-20260314}
\end{table*}

\begin{table*}[ht]
\centering
\scriptsize
\setlength{\tabcolsep}{4pt}
\begin{tabular}{rlrrrrr}
\toprule
N & Horizon & Sparse MLP & Zero MLP & Block-diag.\ MLP & Dense LISTA & Block-diag.\ LISTA \\
\midrule
8 & $H1000$ & 806.715 & 39.2369 & 9.9301 & 323.601 & \textbf{9.6778} \\
8 & $H3000$ & 1.417e+12 & 1.666e+05 & \textbf{25.5180} & 7.644e+10 & 72.2224 \\
16 & $H1000$ & 31.3588 & 9.0779 & \textbf{6.2083} & 13.4294 & 7.1024 \\
16 & $H3000$ & 1.004e+07 & 3695.520 & \textbf{16.8850} & 611.725 & 22.1061 \\
24 & $H1000$ & 7.4924 & 7.2530 & \textbf{5.6117} & 19.0202 & 6.5628 \\
24 & $H3000$ & 7017.227 & 37.9168 & \textbf{19.0987} & 1.323e+04 & 26.1699 \\
32 & $H1000$ & 7.2005 & 6.9907 & \textbf{5.2075} & 91.7743 & 6.1287 \\
32 & $H3000$ & 113.327 & 36.3160 & \textbf{20.4151} & 121.072 & 29.9114 \\
64 & $H1000$ & 208.767 & \textbf{8.1483} & 147.834 & 208.553 & 71.9048 \\
64 & $H3000$ & 224.007 & \textbf{48.9740} & 164.165 & 224.133 & 113.814 \\
\bottomrule
\end{tabular}
\caption{\textbf{Kuramoto dimension sweep at $H1000$ and $H3000$.} Entries are seed-IQM best-periodic errors at $dt=0.00625$ and 200k training, copied from \texttt{results/paper\_seed\_statistics\_20260331/forecasting\_hard\_system\_seed\_stats.csv}. This refreshed table includes the repaired block-diagonal MLP mirror and reports the two horizons most relevant to the moderate- versus long-horizon read. Under IQM, the block-diagonal MLP remains best at $N=16/N24/N32$, while the zero-sparsity MLP is best at $N=64$. Bootstrap 95\% confidence intervals are in the same CSV.}
\label{tab:review-hard-kuramoto-dims-20260314}
\end{table*}

The repaired block-diagonal MLP $N=64$ row was rechecked directly against the
ten raw per-seed evaluation outputs. The catastrophic split
is real rather than a collection typo: seeds 0, 4, 7, and 8 are good
($5.97/24.70$, $5.56/23.92$, $5.89/25.08$, and $6.14/25.30$ at
$H1000/H3000$), while seeds 1, 2, 3, 5, 6, and 9 fail near $208$--$209$ at
$H1000$ and $223$--$225$ at $H3000$. Under the seed-IQM refresh, that split
still leaves the repaired block-diagonal MLP catastrophically worse than the
zero-sparsity MLP at $N=64$.

\begin{table*}[ht]
\centering
\scriptsize
\setlength{\tabcolsep}{3.5pt}
\begin{tabular}{p{0.26\textwidth}p{0.15\textwidth}rrrrrrr}
\toprule
Setting & Model & H100 & H500 & H1000 & H1500 & H2000 & H2500 & H3000 \\
\midrule
\multicolumn{9}{l}{\textit{Hopfield, $N=16$, $P=16$, $dt=0.00625$, 200k}} \\
& Sparse MLP & 0.0437 & \textbf{0.7718} & 2.7977 & 5.0090 & 6.7987 & 8.1932 & 9.3002 \\
& Zero MLP & \textbf{0.0282} & 0.8042 & \textbf{2.6786} & \textbf{4.0959} & \textbf{5.2535} & \textbf{6.2313} & \textbf{7.0298} \\
& Block-diag.\ LISTA & 0.0880 & 1.7772 & 5.9059 & 9.2986 & 11.4717 & 12.8977 & 13.8275 \\
\midrule
\multicolumn{9}{l}{\textit{Hopfield, $N=64$, quarter-step setting $dt=0.0015625$, 200k}} \\
& Sparse MLP & \textbf{0.1775} & \textbf{6.0369} & 70.8941 & 204.301 & 353.973 & 492.496 & 612.102 \\
& Zero MLP & 0.2235 & 6.2662 & \textbf{66.9844} & \textbf{183.587} & \textbf{311.201} & \textbf{428.260} & \textbf{526.925} \\
& Block-diag.\ MLP & 34.4058 & 63.5625 & 218.905 & 417.298 & 669.874 & 850.333 & 996.285 \\
& Dense LISTA & 55.9479 & 91.2265 & 283.748 & 531.207 & 756.366 & 940.576 & 1086.492 \\
& Block-diag.\ LISTA & 39.9454 & 98.2254 & 345.631 & 648.088 & 915.089 & 1131.138 & 1301.535 \\
\bottomrule
\end{tabular}
\caption{\textbf{Hopfield forecasting results.} All entries are seed-IQM best-periodic errors from \texttt{results/paper\_seed\_statistics\_20260331/forecasting\_hard\_system\_seed\_stats.csv}. The $N=16$ setting is the targeted smaller-$dt$ follow-up; the $N=64$ setting is the higher-basin quarter-step probe. The repaired quarter-step block-diagonal MLP control is again strongly negative. Under IQM, the zero-sparsity MLP is best at $H100$ and $H1000$--$H3000$ on $N=16$, while the sparse MLP remains slightly lower only at $H500$; in the quarter-step probe, the sparse MLP is best only at $H100/H500$ and the zero-sparsity MLP takes over from $H1000$ onward. Bootstrap 95\% confidence intervals are in the same CSV.}
\label{tab:review-hard-hopfield-20260314}
\end{table*}


\begin{table*}[ht]
\centering
\scriptsize
\setlength{\tabcolsep}{4pt}
\begin{tabular}{p{0.25\textwidth}rrrrr}
\toprule
Setting & Sparse MLP & Zero MLP & Block-diag.\ MLP & Dense LISTA & Block-diag.\ LISTA \\
\midrule
Kuramoto, $N=16$, $dt=0.00625$, $H1000$ & 30.6451 & 9.1145 & \textbf{6.2799} & 16.8187 & 7.0325 \\
Hopfield, $N=64$, quarter-step, $H1000$ & 70.8941 & \textbf{66.9844} & 218.905 & 283.748 & 345.631 \\
Corrected 4-basin CLV, full 15-seed rebuild, $H1000$ & 0.1520 & \textbf{0.1285} & 0.2891 & 0.1477 & 0.2436 \\
Fixed 8-basin CLV, $dt=0.005$, $H1000$ & 0.5259 & \textbf{0.4027} & 0.5484 & 0.5696 & 9.537e+32 \\
\bottomrule
\end{tabular}
\caption{\textbf{Block-$K$ fairness controls.} These repaired controls isolate whether lowering the error can be explained by imposing block structure in the latent transition alone. Entries are seed-IQM best-periodic errors from \texttt{results/paper\_seed\_statistics\_20260331/forecasting\_hard\_system\_seed\_stats.csv}, with bootstrap 95\% confidence intervals in the same CSV. Under the robust summary, the corrected 4-basin CLV row still flips in favor of the no-sparsity MLP, the quarter-step Hopfield row also favors the no-sparsity MLP, the fixed 8-basin CLV row no longer supports the block-structured controls, and only the Kuramoto row still favors a block-structured control.}
\label{tab:review-hard-controls-20260314}
\end{table*}

\paragraph{Relevant visualizations.}
The most useful figure for this family would be a horizon plot for Kuramoto at
$N=16$ showing sparse MLP, block-diagonal MLP, dense LISTA, and block-diagonal LISTA on the same
axes from $H100$ to $H3000$, because that single plot makes the horizon
dependence clear. A second useful figure would be a compact dimension-sweep
plot for Kuramoto at $H1000$ across $N \in \{8,16,24,32,64\}$. If there is room
for only one additional hard-system visualization, a two-panel figure with
Kuramoto horizon curves on the left and Hopfield horizon curves on the right
would best show that the model ordering changes on Kuramoto but not on
Hopfield.

\subsection{Basin-support and mechanism family}

\paragraph{Aim and isolation logic.}
This family asks whether good forecasting is accompanied by basin-support
alignment: trajectories from the same basin should reuse the same sparse
support, different basins should separate in support space, and recurring
supports should support locally linear latent dynamics. Basin labels are used
only for evaluation, never for training. All models on a given system are
evaluated on the same trajectory corpus, and basin labels are assigned from a
shared long-rollout rule, so differences are attributable to the learned
representation rather than to different sampled trajectories.

\paragraph{Support-alignment audit.}
The labeled support-alignment audit uses 100 trajectories per checkpoint,
trajectory length 500, long-rollout basin assignment after 5000 additional
steps, support threshold $10^{-3}$, and support vectors obtained by averaging the
latent trajectory over time and thresholding the result. In that audit, the
basin mode support is the most common binary support in a basin, mean basin
consistency is the average fraction of trajectories matching that basin mode,
mode uniqueness rate is the fraction of basin pairs with different basin-mode
supports, and cosine separation is mean within-basin cosine similarity minus
mean cosine similarity between basin centroids of the aggregated continuous
latent vectors.

\paragraph{Label-free clustering audit.}
The label-free clustering audit uses 256 trajectories of length 256 and six
explicit feature views: modal support, majority support, last-step support,
last-step cosine, trajectory-mean support, and trajectory-mean cosine. It uses
K-means with the true number of basins, random state 42, 10 restarts, and PCA
to 20 dimensions for continuous feature views when needed.

\paragraph{Recurring-support local-linearity audit.}
The recurring-support local-linearity study uses 256 trajectories of length 256,
majority-support grouping at threshold $10^{-3}$, an 80/20 trajectory split,
minimum retained group sizes of 5 total, 3 train, and 1 test trajectory,
projection to the top 32 principal directions of centered training latent
states, and ridge-regression local linear maps with regularization $10^{-4}$.
It reports
$\mathrm{NRMSE}(y,\hat{y})=\sqrt{\mathrm{mean}((y-\hat{y})^2)}/(\sqrt{\mathrm{mean}(y^2)}+10^{-9})$
for both one-step and 20-step fits, and compares support-conditioned local maps
against a single global map and a shuffled-group baseline that preserves the
retained-group train/test sizes. A retained-coverage gate of 0.60 is used when
deciding whether the result is broad enough for a strong mechanism claim.

\paragraph{Shared trajectory corpus details.}
For the trajectory corpus in the support-alignment and local-linearity studies,
use evaluation seed 42, generate trajectory $i$ from reset seed $42+i$, store
the initial state plus the rollout, and assign the basin label by rolling the
final stored state forward for 5000 additional integration steps. The main
results use random trajectory sampling rather than balanced-per-basin sampling.
The label-free clustering audit also reports a cross-validated
logistic-regression accuracy using up to 5 folds, but never more folds than the
smallest class count in the evaluated basin partition; classes with fewer than
2 examples are dropped only for that classifier score.
The March 21 dependent mechanism extension already applies these same audits to
the matched zero-sparsity MLP baseline. By March 23, the zero-sparsity
support-alignment, label-free clustering, corrected competitive-LV
representation, and recurring-support local-linearity branches are collected,
so the support/mechanism family already contains partial direct
sparse-versus-zero-sparse MLP evidence rather than only an MLP-versus-LISTA
comparison. The repaired direct Kuramoto support-audit rerun is now also
complete, and it stays negative for the zero-sparsity MLP as well.

\paragraph{How to read the numeric summaries.}
The main basin-support read uses four quantities. `ARI' is the adjusted Rand
index from label-free K-means basin recovery, where 1 is perfect and 0 is
chance-like. `CosSep' is mean within-basin cosine similarity minus mean
between-basin cosine similarity on trajectory-mean latent vectors, so positive
values are better. `Cons' is mean basin consistency, namely the fraction of
trajectories in a basin that match that basin's modal support. `Uniq' is the
trajectory-unique support rate, namely the number of distinct trajectory
supports divided by the number of trajectories, so lower is better and 1 means
every trajectory gets its own support. The numeric summary is split so every
data cell stays numeric.

\begin{table*}[ht]
\centering
\scriptsize
\setlength{\tabcolsep}{6pt}
\begin{tabular}{p{0.24\textwidth}p{0.18\textwidth}rr}
\toprule
Case & Model & ARI$_{\min}$ & ARI$_{\max}$ \\
\midrule
\textbf{Multiwell family} & Sparse MLP & 0.953 & 1.000 \\
& Zero MLP & 0.613 & 1.000 \\
& Dense LISTA & 0.752 & 0.990 \\
& Block-diag.\ LISTA & 0.743 & 0.990 \\
\midrule
Duffing & Sparse MLP & 0.247 & 0.247 \\
& Zero MLP & 0.224 & 0.224 \\
& Dense LISTA & 0.247 & 0.247 \\
& Block-diag.\ LISTA & 0.288 & 0.288 \\
\midrule
\textbf{Kuramoto} & Sparse MLP & 0.0017 & 0.0017 \\
& Zero MLP & -0.0001 & -0.0001 \\
& Dense LISTA & 0.0005 & 0.0005 \\
& Block-diag.\ LISTA & 0.0029 & 0.0029 \\
\midrule
\textbf{Corrected competitive LV} & Sparse MLP & 0.105 & 0.105 \\
& Zero MLP & 0.260 & 0.260 \\
& Dense LISTA & 0.083 & 0.083 \\
& Block-diag.\ LISTA & 0.125 & 0.125 \\
\bottomrule
\end{tabular}
\caption{\textbf{Numeric label-free recovery summary.} Higher is better for `ARI'. Multiwell rows report the range across the eight multiwell benchmark systems after taking the best feature-view median for that system and model. Duffing and Kuramoto report the single-system best feature-view median. Corrected competitive Lotka--Volterra reports the best support-view median on the repaired 4-basin follow-up.}
\label{tab:review-mechanism-recovery-20260314}
\end{table*}

\begin{table*}[ht]
\centering
\scriptsize
\setlength{\tabcolsep}{5pt}
\begin{tabular}{p{0.19\textwidth}p{0.16\textwidth}rrrrrr}
\toprule
Case & Model & CosSep$_{\min}$ & CosSep$_{\max}$ & Cons$_{\min}$ & Cons$_{\max}$ & Uniq$_{\min}$ & Uniq$_{\max}$ \\
\midrule
\textbf{Multiwell family} & Sparse MLP & 0.380 & 0.724 & 0.053 & 0.133 & 0.910 & 1.000 \\
& Zero MLP & 0.051 & 0.378 & 0.081 & 0.164 & 0.830 & 0.960 \\
& Dense LISTA & 0.151 & 0.613 & 0.053 & 0.153 & 0.810 & 1.000 \\
& Block-diag.\ LISTA & 0.163 & 0.654 & 0.053 & 0.133 & 0.880 & 1.000 \\
\midrule
Duffing & Sparse MLP & -0.086 & -0.086 & 0.050 & 0.050 & 0.900 & 0.900 \\
& Zero MLP & -0.043 & -0.043 & 0.099 & 0.099 & 0.660 & 0.660 \\
& Dense LISTA & -0.104 & -0.104 & 0.040 & 0.040 & 0.930 & 0.930 \\
& Block-diag.\ LISTA & -0.126 & -0.126 & 0.070 & 0.070 & 0.910 & 0.910 \\
\midrule
\textbf{Kuramoto} & Sparse MLP & -0.300 & -0.300 & 0.424 & 0.424 & 1.000 & 1.000 \\
& Zero MLP & -0.072 & -0.072 & 0.424 & 0.424 & 1.000 & 1.000 \\
& Dense LISTA & -0.259 & -0.259 & 0.424 & 0.424 & 1.000 & 1.000 \\
& Block-diag.\ LISTA & -0.277 & -0.277 & 0.424 & 0.424 & 1.000 & 1.000 \\
\midrule
Hopfield & Sparse MLP & 0.524 & 0.524 & 0.043 & 0.043 & 1.000 & 1.000 \\
& Zero MLP & 0.311 & 0.311 & 0.096 & 0.096 & 0.890 & 0.890 \\
& Dense LISTA & 0.618 & 0.618 & 0.043 & 0.043 & 1.000 & 1.000 \\
& Block-diag.\ LISTA & 0.611 & 0.611 & 0.043 & 0.043 & 1.000 & 1.000 \\
\midrule
\textbf{Corrected competitive LV} & Sparse MLP & -0.056 & -0.056 & 0.107 & 0.107 & 1.000 & 1.000 \\
& Zero MLP & -0.009 & -0.009 & 0.880 & 0.880 & 0.060 & 0.060 \\
& Dense LISTA & -0.061 & -0.061 & 0.107 & 0.107 & 1.000 & 1.000 \\
& Block-diag.\ LISTA & -0.068 & -0.068 & 0.107 & 0.107 & 1.000 & 1.000 \\
\bottomrule
\end{tabular}
\caption{\textbf{Numeric support-reuse summary.} Higher is better for `CosSep' and `Cons'. Lower is better for `Uniq'. Multiwell rows report the range across the eight multiwell benchmark systems. Corrected competitive Lotka--Volterra uses the repaired 4-basin follow-up with the matched zero-sparsity baseline included directly.}
\label{tab:review-mechanism-main-20260314}
\end{table*}

\paragraph{Read by case.}
Multiwell remains the cleanest positive family, but the zero-sparsity baseline
changes \emph{how} it is positive. The sparse MLP keeps the strongest basin
geometry: its best-view recovery stays in the $\mathrm{ARI}=0.953$--$1.000$
range and its cosine separation remains much larger
($0.380$--$0.724$) than the zero-sparsity MLP's
($0.051$--$0.378$). Removing the sparsity penalty does \emph{not} destroy
multiwell recovery, but it does make the supports less cleanly separated even
when they are still often reusable. This is the clearest evidence that explicit
sparsity is acting as a representation-shaping prior rather than as a generic
forecasting stabilizer.

Duffing stays weak. The zero-sparsity MLP improves support reuse relative to
the sparse anchor ($\mathrm{Cons}=0.099$ versus $0.050$ and
$\mathrm{Uniq}=0.66$ versus $0.90$), but its best-view label-free recovery is
still only $\mathrm{ARI}=0.224$, slightly below the sparse anchor's
$0.247$. Here removing the penalty makes supports recur more often without
turning them into cleaner basin identifiers.

Kuramoto is still a hard negative. The zero-sparsity MLP makes cosine
separation less negative ($-0.072$ rather than $-0.300$), but best-view
recovery remains effectively zero and $\mathrm{Uniq}=1.0$ stays saturated. The
completed direct support-audit rerun confirms the same conclusion more
strongly: removing sparsity does not recover reusable non-singleton mode
supports.

Hopfield is the most direct forecasting-versus-mechanism split. The
zero-sparsity MLP forecasts much better than the sparse anchor in the hard
tables and also improves support reuse numerically
($\mathrm{Cons}=0.096$, $\mathrm{Uniq}=0.890$), but it reduces cosine
separation relative to the sparse anchor ($0.311$ versus $0.524$). The sparse
penalty therefore seems helpful for crisp basin geometry on Hopfield even when
it hurts forecasting.

Corrected competitive Lotka--Volterra becomes the one real exception to the old
all-negative MLP story. The zero-sparsity MLP reaches the strongest repaired
4-basin support reuse in this table
($\mathrm{Cons}=0.880$, $\mathrm{Uniq}=0.060$) and the best support-view
recovery score ($\mathrm{ARI}=0.260$), while the sparse anchor and LISTA roots
all remain near the old weak regime. But even here the cosine separation score
stays slightly negative ($-0.009$), so the zero-sparsity model is learning
recurrent supports without producing a fully clean basin-separated continuous
geometry. That is a partial positive, not a universal mechanistic rescue.

\paragraph{Recurring-support local-linearity summary.}
The local-linearity table uses a different set of quantities. `SeedRet' is the
fraction of seeds with at least one retained support group. `Cov' is retained
trajectory coverage. `L20', `G20', and `S20' are the 20-step NRMSE values for
the local, global, and shuffled fits, so lower is better. `W1' is the
fraction of seeds where the one-step local fit beats the matched global fit.
A strong mechanism claim requires $\mathrm{Cov} \ge 0.60$.

\begin{table*}[ht]
\centering
\footnotesize
\setlength{\tabcolsep}{6pt}
\begin{tabular}{lrrrrrr}
\toprule
Model & SeedRet & Cov & L20 & G20 & S20 & W1 \\
\midrule
Sparse MLP & 1.000 & 0.4141 & 0.0205 & 0.0356 & 0.1117 & 1.000 \\
Zero MLP & 1.000 & \textbf{0.4844} & 0.0251 & \textbf{0.0306} & \textbf{0.0649} & 1.000 \\
Block-diagonal LISTA & 1.000 & 0.4336 & \textbf{0.0159} & 0.0380 & 0.0895 & 0.333 \\
Dense LISTA & 1.000 & 0.4414 & 0.0246 & 0.0390 & 0.0784 & 0.333 \\
\bottomrule
\end{tabular}
\caption{\textbf{Recurring-support local-linearity on multiwell strong
transition.} Higher is better for `SeedRet', `Cov', and `W1'. Lower is better
for `L20', `G20', and `S20'.}
\label{tab:review-local-linearity-20260314}
\end{table*}

All four multiwell roots beat their own global and shuffled baselines at 20
steps, but every row stays below the 0.60 coverage gate. The zero-sparsity MLP
actually retains the most trajectories ($\mathrm{Cov}=0.4844$), yet it does
not produce the best local fit: block-diagonal LISTA keeps the best 20-step
error ($\mathrm{L20}=0.0159$), while both MLP roots are the only rows with
$\mathrm{W1}=1.0$. This is another version of the same pattern. Removing the
sparsity penalty makes recurring supports easier to find, but explicit sparsity
still appears to sharpen the local dynamical structure of those supports.
Targeted smaller-step-size Kuramoto is a clean mechanistic negative: 0 of 20
root-seed evaluations retained any support group once the zero-sparsity row is
included, so the local-fit errors are undefined there as well.
